\documentclass[12pt,a4paper]{article}

\usepackage{amsmath,amssymb,amsthm}
\usepackage{mathtools}
\usepackage{enumitem}
\usepackage[margin=2.5cm]{geometry}
\usepackage{hyperref}
\usepackage{booktabs}

% Theorem environments
\newtheorem{theorem}{Theorem}[section]
\newtheorem{proposition}[theorem]{Proposition}
\newtheorem{lemma}[theorem]{Lemma}
\newtheorem{corollary}[theorem]{Corollary}
\theoremstyle{definition}
\newtheorem{definition}[theorem]{Definition}
\newtheorem{example}[theorem]{Example}
\theoremstyle{remark}
\newtheorem{remark}[theorem]{Remark}

% Shortcuts
\newcommand{\C}{\mathbb{C}}
\newcommand{\Z}{\mathbb{Z}}
\newcommand{\F}{\mathbb{F}}
\newcommand{\GL}{\mathrm{GL}}
\newcommand{\End}{\operatorname{End}}
\newcommand{\Hom}{\operatorname{Hom}}
\newcommand{\Ind}{\operatorname{Ind}}
\newcommand{\Stab}{\operatorname{Stab}}
\newcommand{\Hecke}{\mathcal{H}}
\newcommand{\PP}{\mathbb{P}}

\title{The Iwahori--Hecke Algebra\\from Lines in Projective Space}
\author{Notes}
\date{}

\begin{document}
\maketitle

\begin{abstract}
We derive the Iwahori--Hecke algebra $H_n(q)$ from the double coset
algebra $\Hecke(\GL_n(\F_q), B_n(\F_q))$ by computing how the double
coset operators act on the flag variety. The quadratic relation
$(T_s - q)(T_s + 1) = 0$ is derived by counting incidences of lines
in $\PP^1(\F_q)$. The braid relation is derived from the geometry of
flags in $\F_q^3$.
\end{abstract}

\tableofcontents

%% ====================================================================
\section{Setup}
\label{sec:setup}
%% ====================================================================

\subsection{The groups}

Let $q$ be a prime power and $G = \GL_n(\F_q)$ the group of invertible
$n \times n$ matrices over $\F_q$. The \emph{Borel subgroup}
$B = B_n(\F_q)$ is the subgroup of upper triangular matrices:
\[
  B = \left\{\begin{pmatrix} a_{11} & a_{12} & \cdots & a_{1n} \\
  0 & a_{22} & \cdots & a_{2n} \\
  \vdots & & \ddots & \vdots \\
  0 & 0 & \cdots & a_{nn} \end{pmatrix} : a_{ii} \in \F_q^*,\;
  a_{ij} \in \F_q \text{ for } i < j\right\}.
\]
The order is $|B| = (q-1)^n \cdot q^{\binom{n}{2}}$ ($n$ diagonal
entries from $\F_q^*$ and $\binom{n}{2}$ strictly upper-triangular
entries from $\F_q$).

\subsection{The flag variety}

$B$ is the stabiliser of the \emph{standard flag}
\[
  F_\bullet : \quad \{0\} \subset \langle e_1 \rangle \subset
  \langle e_1, e_2 \rangle \subset \cdots \subset
  \langle e_1, \ldots, e_n \rangle = \F_q^n,
\]
where $e_1, \ldots, e_n$ is the standard basis. The coset space
$G/B$ is the \emph{complete flag variety}: the set of all complete
flags
\[
  \{0\} = V_0 \subset V_1 \subset V_2 \subset \cdots \subset V_n =
  \F_q^n, \qquad \dim V_i = i.
\]
For $n = 2$: a complete flag is just a line $V_1 \subset \F_q^2$,
so $G/B = \PP^1(\F_q)$, the projective line with $q + 1$ points.

For $n = 3$: a complete flag is a line $V_1$ inside a plane $V_2$
inside $\F_q^3$, so $G/B$ is the set of incident line-plane pairs.

\subsection{The Weyl group}

The \emph{Weyl group} of $\GL_n$ is $W = S_n$, realised as
permutation matrices. The \emph{simple reflections} are the adjacent
transpositions $s_i = (i,\; i{+}1)$ for $i = 1, \ldots, n-1$,
represented by the permutation matrix that swaps rows $i$ and $i{+}1$.

\subsection{The Bruhat decomposition}

\begin{theorem}[Bruhat]\label{thm:bruhat}
$\GL_n(\F_q)$ decomposes as a disjoint union of double cosets:
\[
  \GL_n(\F_q) = \bigsqcup_{w \in S_n} BwB.
\]
The double cosets are indexed by the Weyl group $W = S_n$.
\end{theorem}

\begin{proof}[Proof for $n = 2$]
$G = \GL_2(\F_q)$, $B$ = upper triangular, $W = \{e, s\}$ where
$s = \bigl(\begin{smallmatrix} 0 & 1 \\ 1 & 0 \end{smallmatrix}\bigr)$.

Any $g = \bigl(\begin{smallmatrix} a & b \\ c & d \end{smallmatrix}\bigr) \in G$:
\begin{itemize}
\item If $c = 0$: then $g$ is upper triangular, so $g \in B = BeB$.
\item If $c \neq 0$: we can write
  $g = \bigl(\begin{smallmatrix} a & b \\ c & d \end{smallmatrix}\bigr)
  = \bigl(\begin{smallmatrix} ac^{-1} & * \\ 1 & * \end{smallmatrix}\bigr)
  \cdot \bigl(\begin{smallmatrix} 0 & 1 \\ 1 & 0 \end{smallmatrix}\bigr)
  \cdot \bigl(\begin{smallmatrix} c & d \\ 0 & * \end{smallmatrix}\bigr)$
  (filling in the $*$'s to make the equation hold). So $g \in BsB$.
\end{itemize}
Thus $\GL_2 = B \sqcup BsB$: every matrix is either upper triangular
or has nonzero $(2,1)$-entry, and in the latter case it belongs to
$BsB$.
\end{proof}

It follows that $\Hecke(\GL_n(\F_q), B)$ has dimension
$|W| = n!$, with basis $\{T_w : w \in S_n\}$ where $T_w$ is the
characteristic function of $BwB$ (suitably normalised).

%% ====================================================================
\section{The quadratic relation from $\PP^1(\F_q)$}
\label{sec:quadratic}
%% ====================================================================

We work with $G = \GL_2(\F_q)$ and $B$ = upper triangular. The Hecke
algebra $\Hecke(G, B)$ is $2$-dimensional with basis $\{T_e, T_s\}$.
We derive the quadratic relation by computing how $T_s$ acts on the
projective line.

\subsection{The action on $\PP^1(\F_q)$}

$\PP^1(\F_q)$ has $q + 1$ points (the lines through the origin in
$\F_q^2$). Let $\delta_L$ denote the basis vector for the point
$L \in \PP^1(\F_q)$ in the permutation representation
$k[\PP^1(\F_q)] = k[G/B]$. The Hecke algebra acts on this space.

\medskip\noindent\textbf{$T_e$} is the identity: $T_e \cdot \delta_L
= \delta_L$.

\medskip\noindent\textbf{$T_s$} acts by sending each point to the sum
of all \emph{other} points:
\begin{equation}\label{eq:Ts-action}
  T_s \cdot \delta_L = \sum_{L' \neq L} \delta_{L'}.
\end{equation}

\begin{proof}[Proof of~\eqref{eq:Ts-action}]
The operator $T_s$ corresponds to the double coset $BsB$. It acts on
$\delta_L$ by summing over all cosets $gB$ with $g^{-1}L' = L$ and
$g \in BsB$, for varying $L'$. Since $B$ stabilises the standard
line $\langle e_1 \rangle$ and $s$ swaps it to $\langle e_2 \rangle$,
the double coset $BsB$ consists of all matrices that send
$\langle e_1 \rangle$ to a line $\neq \langle e_1 \rangle$. By the
transitivity of $B$ on lines $\neq \langle e_1 \rangle$ (upper
triangular matrices act transitively on the $q$ affine points of
$\PP^1$), each line $L' \neq L$ appears exactly once.
\end{proof}

Equivalently: $T_s$ is the adjacency matrix of the complete graph
$K_{q+1}$ on $\PP^1(\F_q)$, minus the identity. In matrix form,
$T_s = J - I$ where $J$ is the all-ones matrix and $I$ the identity.

\subsection{Deriving the quadratic relation}

\begin{theorem}\label{thm:quadratic}
$T_s^2 = (q-1)\, T_s + q\, T_e$.
\end{theorem}

\begin{proof}
Apply $T_s$ twice to $\delta_L$ for any point $L \in \PP^1(\F_q)$.

\medskip\noindent\textbf{First application}: by~\eqref{eq:Ts-action},
$T_s \cdot \delta_L = \sum_{L' \neq L} \delta_{L'}$ (a sum of $q$
terms).

\medskip\noindent\textbf{Second application}: apply $T_s$ to each
$\delta_{L'}$:
\[
  T_s^2 \cdot \delta_L
  = \sum_{L' \neq L} T_s \cdot \delta_{L'}
  = \sum_{L' \neq L}\; \sum_{L'' \neq L'} \delta_{L''}.
\]
Now count how many times each point $L''$ appears in this double sum:

\begin{itemize}
\item \textbf{$L'' = L$}: for each $L' \neq L$, the inner sum
  includes $L'' = L$ (since $L \neq L'$). There are $q$ such $L'$.
  Contribution: $q \cdot \delta_L$.

\item \textbf{$L'' \neq L$}: for a fixed $L'' \neq L$, the inner sum
  includes $L''$ whenever $L' \neq L''$. Among the $q$ choices of
  $L'$ (all $\neq L$), exactly one is excluded ($L' = L''$), leaving
  $q - 1$ choices.
  Contribution: $(q-1) \cdot \delta_{L''}$ for each $L'' \neq L$.
\end{itemize}

Combining:
\[
  T_s^2 \cdot \delta_L
  = q \cdot \delta_L + (q-1) \sum_{L'' \neq L} \delta_{L''}
  = q \cdot T_e \cdot \delta_L + (q-1) \cdot T_s \cdot \delta_L.
\]
Since this holds for every $L$, we have $T_s^2 = (q-1)\,T_s + q\,T_e$.
\end{proof}

\begin{corollary}\label{cor:quadratic-factored}
The quadratic relation factors as
\begin{equation}\label{eq:quadratic}
  (T_s - q)(T_s + 1) = 0.
\end{equation}
\end{corollary}

\begin{proof}
$T_s^2 - (q-1)T_s - q = T_s^2 - qT_s + T_s - q
= T_s(T_s - q) + (T_s - q) = (T_s + 1)(T_s - q) = 0$.
\end{proof}

\begin{remark}
The two roots $T_s = q$ and $T_s = -1$ are the eigenvalues of
$T_s$ on $k[\PP^1(\F_q)]$:
\begin{itemize}
\item $T_s = q$ on the trivial representation (the constant function
  $\sum_L \delta_L$): $T_s$ sends it to $q$ times itself, since each
  point has $q$ neighbours.
\item $T_s = -1$ on the $q$-dimensional complement (functions summing
  to zero): $T_s \cdot f = (J - I) \cdot f = -f$ since $J \cdot f = 0$
  for $f$ orthogonal to the constant function.
\end{itemize}
\end{remark}

\subsection{The limit $q \to 1$}

At $q = 1$ (formally): the relation becomes $(T_s - 1)(T_s + 1) = 0$,
i.e., $T_s^2 = 1$. This is the involution relation for the symmetric
group. The projective line $\PP^1(\F_1)$ has $1 + 1 = 2$ points, and
$T_s$ swaps them---exactly the transposition $(1\;2) \in S_2$.

%% ====================================================================
\section{The braid relation from flags in $\F_q^3$}
\label{sec:braid}
%% ====================================================================

For $G = \GL_3(\F_q)$, the Weyl group is $W = S_3$ with generators
$s_1 = (1\;2)$ and $s_2 = (2\;3)$. The Hecke algebra
$\Hecke(G, B)$ has dimension $6$ with basis $\{T_w : w \in S_3\}$.

We need to verify: $T_{s_1} T_{s_2} T_{s_1} = T_{s_2} T_{s_1}
T_{s_2}$ (the braid relation).

\subsection{The flag variety for $\GL_3$}

A complete flag in $\F_q^3$ is an incident pair (line $\ell$,
plane $\pi$) with $\ell \subset \pi \subset \F_q^3$. The standard
flag is $F_\bullet = (\langle e_1 \rangle,\;
\langle e_1, e_2 \rangle)$.

The operators $T_{s_1}$ and $T_{s_2}$ act on flags by modifying one
component:
\begin{itemize}
\item $T_{s_1}$ changes the \emph{line} $\ell$ while keeping the
  plane $\pi$ fixed: it sums over all lines $\ell' \neq \ell$ in
  $\pi$.
\item $T_{s_2}$ changes the \emph{plane} $\pi$ while keeping the
  line $\ell$ fixed: it sums over all planes $\pi' \neq \pi$
  containing $\ell$.
\end{itemize}

\begin{proof}[Why $T_{s_1}$ changes only the line]
The simple reflection $s_1$ swaps $e_1$ and $e_2$. So $BsB$ consists
of matrices that move $\langle e_1 \rangle$ to another line inside
$\langle e_1, e_2 \rangle$, while preserving $\langle e_1, e_2
\rangle$ itself. That is, $T_{s_1}$ replaces the line component of the
flag by another line in the same plane.
\end{proof}

\subsection{Counting}

In $\F_q^3$:
\begin{itemize}
\item A plane $\pi \cong \F_q^2$ contains $q + 1$ lines. So
  $T_{s_1}$ sums over $q$ lines (all lines in $\pi$ except $\ell$).
\item Through a line $\ell$, there are $q + 1$ planes in $\F_q^3$.
  So $T_{s_2}$ sums over $q$ planes (all planes through $\ell$
  except $\pi$).
\end{itemize}

Each operator is a ``local $T_s$'': $T_{s_1}$ acts like the
$\PP^1(\F_q)$-operator from Section~\ref{sec:quadratic} on the
lines within a fixed plane, and $T_{s_2}$ acts like it on the planes
through a fixed line. Each satisfies the quadratic relation
$(T_{s_i} - q)(T_{s_i} + 1) = 0$ by the same counting argument as
Theorem~\ref{thm:quadratic}.

\subsection{The braid relation}

\begin{theorem}\label{thm:braid}
$T_{s_1} T_{s_2} T_{s_1} = T_{s_2} T_{s_1} T_{s_2}$.
\end{theorem}

\begin{proof}
Apply both sides to a flag $(\ell, \pi)$ and show they give the same
result.

\medskip\noindent\textbf{Left side: $T_{s_1} T_{s_2} T_{s_1}$.}
Starting from $(\ell, \pi)$:
\begin{enumerate}
\item $T_{s_1}$: replace $\ell$ by another line $\ell'$ in $\pi$
  (sum over $q$ choices).
\item $T_{s_2}$: replace $\pi$ by another plane $\pi'$ through
  $\ell'$ (sum over $q$ choices).
\item $T_{s_1}$: replace $\ell'$ by another line $\ell''$ in $\pi'$
  (sum over $q$ choices).
\end{enumerate}
Result: sum over all flags $(\ell'', \pi')$ reachable by a path
$(\ell, \pi) \to (\ell', \pi) \to (\ell', \pi') \to (\ell'', \pi')$,
where each step changes one component.

\medskip\noindent\textbf{Right side: $T_{s_2} T_{s_1} T_{s_2}$.}
Starting from $(\ell, \pi)$:
\begin{enumerate}
\item $T_{s_2}$: replace $\pi$ by $\pi'$ through $\ell$ ($q$ choices).
\item $T_{s_1}$: replace $\ell$ by $\ell'$ in $\pi'$ ($q$ choices).
\item $T_{s_2}$: replace $\pi'$ by $\pi''$ through $\ell'$ ($q$ choices).
\end{enumerate}
Result: sum over all flags $(\ell', \pi'')$ reachable by a path
$(\ell, \pi) \to (\ell, \pi') \to (\ell', \pi') \to (\ell', \pi'')$.

\medskip\noindent Both sides produce a sum over flags at
``Bruhat distance 3'' from $(\ell, \pi)$. The key point: a flag
$(\ell'', \pi')$ is at distance 3 if and only if $\ell'' \neq \ell$
and $\pi' \neq \pi$ and $\ell'' \not\subset \pi$ and
$\ell \not\subset \pi'$. This set---and the multiplicity with which
each flag appears---is the same for both orderings, because the
composition $s_1 s_2 s_1 = s_2 s_1 s_2$ in $S_3$ means the two
paths through the Bruhat graph reach the same double coset $Bs_1 s_2
s_1 B = B w_0 B$ (where $w_0$ is the longest element of $S_3$).

More precisely: both $T_{s_1}T_{s_2}T_{s_1}$ and $T_{s_2}T_{s_1}T_{s_2}$
equal $T_{w_0}$ (the basis element for the longest Weyl group element),
because in the Hecke algebra one has $T_{w} = T_{s_{i_1}} \cdots
T_{s_{i_k}}$ whenever $w = s_{i_1} \cdots s_{i_k}$ is a reduced
expression, and $w_0 = s_1 s_2 s_1 = s_2 s_1 s_2$ are the two
reduced expressions.
\end{proof}

\begin{remark}
The proof relies on the fact that $\ell(s_1 s_2 s_1) = 3 =
\ell(s_1) + \ell(s_2) + \ell(s_1)$, where $\ell(w)$ is the length
of $w$ (number of simple reflections in a reduced expression). This
ensures that $T_{s_1}T_{s_2}T_{s_1}$ is a single basis element
$T_{w_0}$, not a sum of basis elements. The multiplicativity
$T_u T_v = T_{uv}$ holds whenever $\ell(uv) = \ell(u) + \ell(v)$.
\end{remark}

%% ====================================================================
\section{The isomorphism}
\label{sec:iso}
%% ====================================================================

\subsection{The Iwahori--Hecke algebra by generators and relations}

\begin{definition}
The \emph{Iwahori--Hecke algebra} $H_n(q)$ is the associative algebra
over $k$ generated by elements $T_1, \ldots, T_{n-1}$ subject to:
\begin{align}
  (T_i - q)(T_i + 1) &= 0
  && \text{for all } i, \label{eq:IH-quad} \\
  T_i T_{i+1} T_i &= T_{i+1} T_i T_{i+1}
  && \text{for all } i, \label{eq:IH-braid} \\
  T_i T_j &= T_j T_i
  && \text{for } |i - j| > 1. \label{eq:IH-comm}
\end{align}
\end{definition}

At $q = 1$: the quadratic relation becomes $T_i^2 = 1$ and the
algebra is $\C[S_n]$.

\subsection{The isomorphism theorem}

\begin{theorem}[Iwahori, 1964]\label{thm:iwahori}
There is an algebra isomorphism
\[
  H_n(q) \;\xrightarrow{\;\sim\;}\;
  \Hecke(\GL_n(\F_q),\; B_n(\F_q)),
\]
sending the generator $T_i$ to the (suitably normalised) characteristic
function of the double coset $B s_i B$.
\end{theorem}

\begin{proof}[Proof sketch]
We must verify that the double coset operators $T_{s_i}$ satisfy
the three relations of $H_n(q)$.

\begin{enumerate}
\item \emph{Quadratic relation}~\eqref{eq:IH-quad}: proved in
  Theorem~\ref{thm:quadratic} for $n = 2$. The same argument works
  for any $i$: the operator $T_{s_i}$ acts on the $i$-th component
  of the flag (a line in a plane), and the counting on $\PP^1(\F_q)$
  gives $(T_{s_i} - q)(T_{s_i} + 1) = 0$.

\item \emph{Braid relation}~\eqref{eq:IH-braid}: proved in
  Theorem~\ref{thm:braid} for $(s_1, s_2)$ in $\GL_3$. The same
  argument works for any adjacent pair $(s_i, s_{i+1})$: both sides
  equal $T_{s_i s_{i+1} s_i}$ because the product is length-additive.

\item \emph{Commutativity}~\eqref{eq:IH-comm}: for $|i - j| > 1$,
  the operators $T_{s_i}$ and $T_{s_j}$ act on disjoint components
  of the flag (non-adjacent levels), so they commute.
\end{enumerate}

It remains to show the map is an isomorphism. Both algebras have
dimension $n!$: the Iwahori--Hecke algebra by its standard basis
$\{T_w : w \in S_n\}$ (defined by $T_w = T_{s_{i_1}} \cdots
T_{s_{i_k}}$ for any reduced expression), and the Hecke algebra by
the Bruhat decomposition (Theorem~\ref{thm:bruhat}). The map sends
basis to basis.
\end{proof}

%% ====================================================================
\section{Summary}
\label{sec:summary}
%% ====================================================================

The Iwahori--Hecke algebra $H_n(q)$ has two faces:

\begin{center}
\renewcommand{\arraystretch}{1.5}
\begin{tabular}{@{}p{5.5cm}p{5.5cm}@{}}
\toprule
\textbf{Abstract (generators \& relations)}
  & \textbf{Geometric (double cosets)} \\
\midrule
Generators $T_1, \ldots, T_{n-1}$
  & Double cosets $Bs_iB$ \\
Quadratic $(T_i - q)(T_i + 1) = 0$
  & Counting on $\PP^1(\F_q)$: $q$ neighbours \\
Braid $T_iT_{i+1}T_i = T_{i+1}T_iT_{i+1}$
  & Flag geometry: two paths to $w_0$ \\
Commuting $T_iT_j = T_jT_i$ for $|i-j|>1$
  & Non-adjacent flag components \\
$q = 1$ gives $\C[S_n]$
  & $\F_1$: Borel = trivial, flags = permutations \\
\bottomrule
\end{tabular}
\end{center}

The parameter $q$ is the number of elements in $\F_q$, which is the
number of points in a ``local $\PP^1$'' (the lines in a plane, or the
planes through a line). The quadratic relation $(T_s - q)(T_s + 1) = 0$
encodes this: $T_s$ acting on $\PP^1(\F_q)$ has eigenvalue $q$ on the
constant function and eigenvalue $-1$ on the complement. As $q \to 1$,
$\PP^1$ shrinks to two points, $T_s$ becomes a transposition, and
$H_n(q)$ degenerates to $\C[S_n]$.

\begin{thebibliography}{99}

\bibitem{Bourbaki}
N.~Bourbaki,
\emph{Lie Groups and Lie Algebras}, Chapters 4--6,
Springer-Verlag, 2002.

\bibitem{Humphreys}
J.~E.~Humphreys,
\emph{Reflection Groups and Coxeter Groups},
Cambridge Studies in Advanced Mathematics 29,
Cambridge University Press, 1990.

\bibitem{Iwahori}
N.~Iwahori,
\emph{On the structure of a Hecke ring of a Chevalley group over a
finite field},
J.~Fac.\ Sci.\ Univ.\ Tokyo Sect.\ I \textbf{10} (1964), 215--236.

\end{thebibliography}

\end{document}
